% A worked example: the orbifold C/Z_5 and the section z -> z^2.
% Compile with xelatex (twice, for the references).
\documentclass[11pt,a4paper]{article}

\usepackage[margin=2.6cm]{geometry}
\usepackage{amsmath,amssymb,amsthm}
\usepackage{unicode-math}
\setmainfont{STIXTwoText}[Extension=.otf,UprightFont=*-Regular,BoldFont=*-Bold,
  ItalicFont=*-Italic,BoldItalicFont=*-BoldItalic]
\setmathfont{STIXTwoMath-Regular.otf}
\setmonofont[Scale=0.84]{Menlo}
\usepackage{xcolor}
\usepackage{url}
\usepackage{array,longtable,booktabs}
\usepackage{enumitem}
\usepackage[hidelinks]{hyperref}

\definecolor{leangrey}{gray}{0.40}

% Lean declaration names: monospace, line breaks after dots and underscores.
\DeclareUrlCommand\lean{\def\UrlFont{\ttfamily\upshape}}
\renewcommand{\UrlBreaks}{\do\.\do\_}
\renewcommand{\UrlBigBreaks}{}
% A small grey reference to the Lean declaration, set flush right at the end of a statement.
\newcommand{\leanref}[1]{\par\nopagebreak
  {\raggedleft\small\upshape\color{leangrey}Lean:~#1\par}}

\setlength{\emergencystretch}{3em}
\setlist{itemsep=2pt,topsep=4pt}

\theoremstyle{plain}
\newtheorem{theorem}{Theorem}[section]
\newtheorem{proposition}[theorem]{Proposition}
\newtheorem{lemma}[theorem]{Lemma}
\theoremstyle{definition}
\newtheorem{definition}[theorem]{Definition}
\theoremstyle{remark}
\newtheorem{remark}[theorem]{Remark}

\newcommand{\C}{\mathbb{C}}
\newcommand{\R}{\mathbb{R}}
\newcommand{\Q}{\mathbb{Q}}
\newcommand{\Z}{\mathbb{Z}}
\newcommand{\N}{\mathbb{N}}
\newcommand{\ZZ}{\Z/5}
\newcommand{\abs}[1]{\lvert #1\rvert}
\newcommand{\card}{\sharp}

\title{A worked example: the orbifold $\C/\Z_5$ and the section $z\mapsto z^2$}
\author{}
\date{A note accompanying a Lean formalization of\\
H.~Hofer, K.~Wysocki, E.~Zehnder, \emph{Polyfold and Fredholm Theory}, Springer, 2021}

\begin{document}
\maketitle

\noindent
This note carries out Theorem 16.6.6 (p.~739) of \emph{Polyfold and Fredholm Theory} for Fredholm
index $0$ on one small polyfold: the orbit space $\C/\Z_5$ of the action of the fifth roots of unity,
with the bundle of weight $2$ and the section $z\mapsto z^2$. Every step is made explicit, from the
ep-groupoid and its strong bundle, through the hypotheses of the theorem, the choice of an admissible
sc$^+$-multisection and the canonical orientation of the perturbed solution set, to the value of the
invariant, which is $2/5$. The example is small, but it is a genuine orbifold case: no single-valued
equivariant perturbation is transversal, so the multisections of the book are needed. A second version
adds an infinite-dimensional factor with strictly nested levels, so that sc-calculus and sc-Fredholm
theory enter in earnest; the same conclusions hold there. All statements below are machine-checked
in Lean~4 with Mathlib, without unproved assumptions and with only the standard axioms
(propositional extensionality, choice, quotient soundness). Book references use the numbering and
pages of the 2021 edition. Lean names refer to the namespace
\lean{Polyfold.Examples.OrbifoldSquare} unless written in full; the files are in
\lean{Polyfold/Examples/}, and the appendix lists them.

%======================================================================
\section{The example}

Let $\zeta=e^{2\pi i/5}$ and let $G=\ZZ$ act on $\C$ by $j\cdot z=\zeta^j z$.

\paragraph{The ep-groupoid.}
$X=G\ltimes\C$ is the action groupoid: its objects are the points of $\C$, its morphisms are the pairs
$(j,z)\in G\times\C$, with source $z$ and target $\zeta^j z$. The sc-structure on $\C$ is the constant
one, $\C_m=\C$ for all $m$. The groupoid $X$ is étale and proper; it is effective, the isotropy group of
$0$ is $G_0\cong\ZZ$, and every $z\neq0$ has trivial isotropy. Its orbit space $Z=\abs{X}=\C/\Z_5$ is a
polyfold without boundary.

\paragraph{The strong bundle.}
$W=\C\times\C\to\C$, $(z,w)\mapsto z$, with the constant double filtration $W_{m,k}=\C\times\C$ and the
action of weight $2$,
\[
  j\cdot(z,w)=(\zeta^j z,\ \zeta^{2j}w).
\]
It descends to a strong polyfold bundle $p\colon Y=\abs{W}\to Z$.

\paragraph{The section.}
$f(z)=(z,z^2)$. It is equivariant, since $(\zeta^j z)^2=\zeta^{2j}z^2$, and so defines a section of
$p$. Its linearization at $z$ is the real-linear map $h\mapsto 2zh$ of $\C=\R^2$; hence $f$ is
sc-Fredholm with tame germ data, of index $2-2=0$ at every point. The solution set is
$S_f=\{0\}$, which is compact.

\paragraph{Auxiliary norm and orientation.}
The auxiliary norm is $N(z,w)=\abs{w}$ on $W_{0,1}$, which is invariant under the action, and descends to
an auxiliary norm $n$ on $Y$. We take $U=Z$. For $y\in\C$, the linearization at $y$ of $f-s$, where $s$
is the constant sc$^+$-section with $s(y)=f(y)$, is multiplication by $2y$; for $y\neq0$ it is a
complex-linear isomorphism of $\C$. For an isomorphism $T$ the determinant line is canonically $\R$, and
an orientation of it is \emph{positive} if it is the orientation of $1$. We write $o$ for the continuous
orientation of $f$ (Definition 16.6.5, p.~739) that is positive at these complex-linear isomorphisms.

%======================================================================
\section{Why multisections are needed}

A single-valued sc$^+$-section of $W$ compatible with the morphisms has the form $s(z)=(z,\psi(z))$ with
$\psi$ of class $C^1$ and $\psi(\zeta z)=\zeta^2\psi(z)$. Its multisection is the characteristic function
of the graph of $s$.

\begin{proposition}\label{prop:obstruction}
For every such $s$ we have $s(0)=f(0)$, so $0$ remains a solution of the perturbed problem, and
$(f-s)'(0)$ is not onto. In particular the multisection of $s$ is not in general position with $f$
at $0$.
\leanref{\lean{not_isTransversalAt_sectionMulti}, \lean{linear_eq_zero_of_equivariant}}
\end{proposition}

\begin{proof}
At $z=0$ the relation $\psi(\zeta z)=\zeta^2\psi(z)$ gives $\psi(0)=\zeta^2\psi(0)$; since
$\zeta^2\neq1$, $\psi(0)=0$, and $f(0)=(0,0)=s(0)$. Since $f'(0)=0$, the linearization is
$(f-s)'(0)=-D\psi(0)=:L$, a real-linear map of $\C$ satisfying $L(\zeta h)=\zeta^2L(h)$. Write
$L(h)=ah+b\bar h$ with $a,b\in\C$. Then
\[
  L(\zeta h)=a\zeta h+b\zeta^{-1}\bar h,\qquad \zeta^2L(h)=a\zeta^2h+b\zeta^2\bar h,
\]
and comparing coefficients, $a(\zeta-\zeta^2)=0$ and $b(\zeta^{-1}-\zeta^2)=0$. As
$2\not\equiv\pm1\pmod 5$, both factors are non-zero, so $a=b=0$ and $L=0$, which is not onto $\C$.
\end{proof}

The obstruction disappears for multisections. For $c\in\C$, $c\neq0$, let
\[
  \Lambda_c(z,w)=\tfrac15\,\card\{\,j\in\ZZ : w=\zeta^{2j}c\,\},
\]
the multisection with the five constant branches $s_j(z)=(z,\zeta^{2j}c)$, $j\in\ZZ$, each of weight
$1/5$. Since $2$ is prime to $5$, the values $\zeta^{2j}c$ are the five points of $\mu_5c$, and the
action permutes the branches, $j\cdot s_k = s_{j+k}\circ(\zeta^j\,\cdot)$. Thus $\Lambda_c$ is an
sc$^+$-multisection, and it descends to an sc$^+$-multisection $\lambda_c$ of $p$ (Definition 16.5.1,
p.~733).

%======================================================================
\section{What is proved}

An sc$^+$-multisection $\lambda$ of $p$ is called \emph{admissible} if it is controlled by $(n,U)$
(Definition 16.5.4, p.~734) and $(f,\lambda)$ is in general position (Definition 15.3.5, p.~690), as in
Theorem 16.6.6. The perturbed solution set $\theta=\lambda\circ f$ then has dimension $0$; with the
canonical orientation $\hat\theta$ inherited from $(f,o)$ (Theorem 15.4.1 and Definition 15.4.2,
pp.~701--702), the invariant of index $0$ is the branched integral $\int_{\hat\theta}\tau$ of a closed
$0$-form $\tau$ (Definition 9.5.6, p.~439, for $n=0$), which we call the \emph{count}.

\begin{theorem}\label{thm:hyp}
The hypotheses of Theorem 16.6.6 for index $0$ hold for $(p,f,o,n,U)$: $\partial Z=\emptyset$, the
solution set of $f$ is compact, $n$ is an auxiliary norm, $U=Z$ is open, $(n,U)$ controls compactness,
and the orientation $o$ is given on a model with tame germ data and index $0$. Consequently the
conclusion of Theorem 16.6.6 holds for this polyfold bundle and section.
\leanref{\lean{theorem_16_6_6_hypotheses}, \lean{theorem_16_6_6_sq}}
\end{theorem}

\begin{proof}[Sketch]
$\C$ has no boundary and $S_f=\{0\}$. The norm $N$ is continuous on $W_{0,1}$, restricts to a norm on
each fibre and is invariant. For compactness control: $N(f(z))\le1$ means $\abs{z}^2\le1$, a compact
set. The model of $o$ is the graph section $z\mapsto(z,z^2)$ of a smooth map between finite-dimensional
spaces, which is sc-Fredholm of index $0$ with tame germ data.
\end{proof}

\begin{theorem}\label{thm:admissible}
For $0<\abs c\le1$ the multisection $\lambda_c$ is admissible, and its perturbed solution set carries
the canonical orientation of $(f,o)$. Any two admissible perturbations with the canonical orientation
have the same count of a closed $0$-form $\tau$; that common value is the count over $\lambda_c$, for
every $0<\abs c\le1$.
\leanref{\lean{exists_canonical_t5}, \lean{count_eq_of_isAdmissible}, \lean{exists_unique_count}}
\end{theorem}

\begin{proof}[Sketch]
Every branch has norm $N(s_j(z))=\abs c\le1$, and the domain support lies in $U=Z$, so $\lambda_c$ is
controlled by $(n,Z)$. A perturbed solution is a $y$ with $y^2=\zeta^{2j}c$ for some $j$; then
$y\neq0$ and $(f-s_j)'(y)$ is multiplication by $2y$, which is onto; there is no boundary. Hence
$(f,\lambda_c)$ is in general position, and Theorem 15.4.1 provides the canonical orientation. The
independence statement is Theorem 16.6.6 for index $0$, applied through Theorem~\ref{thm:hyp}.
\end{proof}

\begin{theorem}\label{thm:count}
For every admissible perturbation with the canonical orientation of $(f,o)$, the count of the constant
form $1$ is
\[
  \int_{\hat\theta}1=\frac25 .
\]
\leanref{\lean{count_eq_two_fifths}, \lean{count_eq_two_fifths_all}}
\end{theorem}

\begin{proof}
By Theorem~\ref{thm:admissible} it suffices to compute with $\lambda_1$, and the constant form $1$ is
closed (Lemma~\ref{lem:const}). The perturbed solutions are the $y$ with $y^2=\zeta^{2j}$, that is
$y=\pm\zeta^j$: ten points, forming the two orbits of $1$ and of $-1$, which are different points of
$Z$ because $-1$ is not a fifth root of unity. Each of them has trivial isotropy, being non-zero. Each
lies on exactly one branch: $y^2$ determines $j$, as the $\zeta^{2j}$ are distinct. So the weight at
each solution is $\Lambda_1(f(y))=1/5$.

At such a point the sheet has dimension $0$, and by Lemma~\ref{lem:sign} the sign of the canonical
orientation is $+1$ exactly when $o$ is positive at $(f-s_j)'(y)$. Since $s_j$ is constant, this is
multiplication by $2y$, a complex-linear isomorphism, at which $o$ is positive by construction. The
signed weight at each solution is therefore $+1/5$. The formula for a finite orbit support
(Lemma~\ref{lem:finite}), with one representative of each orbit, gives
\[
  \int_{\hat\theta}1=\frac1{\card G_1}\cdot\frac15\cdot1+\frac1{\card G_{-1}}\cdot\frac15\cdot1
  =\frac15+\frac15=\frac25 .\qedhere
\]
\end{proof}

\begin{remark}
The value agrees with the local degree $2$ of $z\mapsto z^2$ at $0$ divided by the order $5$ of the
isotropy group $G_0$. This description of the answer is a remark; it is not one of the formal
statements.
\end{remark}

\begin{remark}
The existence of $o$ rests on Lemma~\ref{lem:path}: any two points $y_0,y_1\neq0$ are joined by a path
$t\mapsto y_0e^{tw}$ in $\C\setminus\{0\}$, along which the linearizations $2y_0e^{tw}$ are isomorphisms,
so positivity at one of them is carried to all. The formal development
records $o$, its model and the positivity at all perturbed solutions
(\lean{sqOrientationPos}, \lean{sqOrientationRepPos}, \lean{positiveZeros_sqOrientationRepPos}).
\end{remark}

%======================================================================
\section{The infinite-dimensional version}

To let the sc-calculus of Part~I, the shift of the strong bundle and sc-Fredholm theory in infinite
dimensions enter, the example is crossed with an sc-Hilbert space whose levels are strictly nested
(namespace \lean{Polyfold.Examples.OrbifoldSquareInf}).

\paragraph{The sc-Hilbert space.}
Let $H$ be the space of real sequences $x=(x_k)_{k\ge0}$ with the levels
\[
  H_m=\Bigl\{x:\ \abs{x}_m^2=\sum_{k\ge0}(k+1)^{2m}x_k^2<\infty\Bigr\}.
\]
The inclusions $H_{m+1}\to H_m$ are compact (the weights $(k+1)^{-1}$ tend to $0$) and not onto; the
smooth points $H_\infty=\bigcap_m H_m$ are dense in every level. The sequence $x_k=1/(k+1)$ lies in $H_0$
and not in $H_1$, so not every object is smooth.
\leanref{\lean{Polyfold.Examples.WeightedEllTwo.H},
\lean{Polyfold.Examples.WeightedEllTwo.not_surjective_inclLevel},
\lean{Polyfold.Examples.WeightedEllTwo.harmH_not_mem_smooth}, \lean{exists_not_mem_smoothSet}}

\paragraph{The polyfold bundle.}
The ep-groupoid is $(G\ltimes\C)\times H$, with $G$ acting trivially on $H$. The strong bundle has fibre
$\C\times H$, of weight $2$ on $\C$ and trivial on $H$, with the double filtration
$(\C\times H)_{m,k}=\C\times H_k$ over level $m$, $0\le k\le m+1$. The section is
\[
  F(z,e)=\bigl((z,e),(z^2,e)\bigr).
\]
Its $H$-part is the identity of $H$, which is sc-Fredholm with tame germ data (it is conjugate to a
basic germ), so $F$ is sc-Fredholm of index $0$. The auxiliary norm is
$N\bigl((z,e),(w,h)\bigr)=\max(\abs w,\abs h_1)$, built from the $H_1$-norm of the fibre; its unit ball
is precompact in $H_0$ because the inclusion $H_1\to H_0$ is compact, and this is what gives compactness
control. The orientation is the continuous one positive at the isomorphisms $(2z,\mathrm{id})$,
$z\ne0$.
\leanref{\lean{infBundle}, \lean{isScFredholmSectionTame_idSec}, \lean{isConjugateToBasicGerm_id},
\lean{isIndexAt_sqInf}}

\begin{theorem}\label{thm:inf}
For the bundle $(G\ltimes\C)\times H$ and the section $F$:
\begin{enumerate}[label=(\roman*)]
\item every hypothesis of Theorem 16.6.6 for index $0$ holds as stated, and so does its conclusion;
\item for $0<\abs c\le1$ the five-branch multisection with branches $(\zeta^{2j}c,0)$, $j\in\ZZ$, of
  weight $1/5$, is admissible and carries the canonical orientation;
\item every admissible perturbation with the canonical orientation has count
  $\int_{\hat\theta}1=2/5$.
\end{enumerate}
\leanref{\lean{theorem_16_6_6_hypotheses_inf}, \lean{theorem_16_6_6_inf},
\lean{isInGeneralPosition_Λinf}, \lean{exists_canonical_tinf}, \lean{count_eq_two_fifths_inf},
\lean{count_eq_two_fifths_inf_all}}
\end{theorem}

\begin{proof}[Sketch]
The perturbed solutions are the $(y,0)$ with $y^2=\zeta^{2j}c$; at each the linearization is
$(2y,\mathrm{id})$, an isomorphism of $\C\times H$. The count reduces to the finite-dimensional one:
two orbits, trivial isotropy, weight $1/5$, sign $+1$.
\end{proof}

The proposition of \S2 is stated for the finite-dimensional example. In the product an equivariant
sc$^+$-section need not vanish on the $H$-component; the obstruction sits on the $\C$-component.

%======================================================================
\section{General lemmas obtained on the way}

The following hold for any strong bundle over an ep-groupoid, or any polyfold, and are used above.

\begin{lemma}[Counts over finitely many orbits]\label{lem:finite}
Let $\hat\Theta=(\Theta,\hat T)$ be an oriented, tame, compact branched ep$^+$-subgroupoid of
dimension $0$, and let $s$ be a finite set of points of $\operatorname{supp}\Theta$ with pairwise
different orbits whose orbits cover the orbit support. Then for every $0$-form $\omega$,
\[
  \int_{\hat\Theta}\omega=\sum_{x\in s}\frac1{\card G_x}\,\varepsilon\bigl(\hat T(x)\bigr)\,\omega(x),
\]
where $\varepsilon(\hat T(x))$ is the signed weight of $\hat T(x)$: the sum of the weights of the
branches through $x$, each with the sign of its orientation. This is Definition 9.5.6 (p.~439) for
$n=0$.
\leanref{\lean{Polyfold.EP.pointBranchedIntegral_eq_sum},
\lean{Polyfold.SBLocalization.PointSolution.count_eq_sum}}
\end{lemma}

\begin{lemma}[$d$ of constants]\label{lem:const}
On every M-polyfold, the constant $0$-forms are closed: $d\,c=0$. The same holds on ep-groupoids.
\leanref{\lean{Polyfold.MPolyfold.extD_constForm}, \lean{Polyfold.EP.epFormsInftyD_constFormInfty}}
\end{lemma}

\begin{proof}
A constant form is the pull-back of a form on the point $\R^0$ by the constant map, $d$ commutes with
pull-back, and every $1$-form on $\R^0$ vanishes.
\end{proof}

\begin{lemma}[The sign in dimension $0$]\label{lem:sign}
In the setting of Theorem 15.4.1 (pp.~701--702), let $y$ lie on a sheet $M_i$ of dimension $0$. The
sign of the oriented tangent space $(T_yM_i,o_{i,y})$ is $+1$ if and only if $o_y$ is positive at the
isomorphism $(f-s_i)'(y)$. If $o_y$ is positive at every branch through $y$, then the signed weight of
$\hat T(y)$ equals $\Lambda(f(y))$.
\leanref{\lean{Polyfold.SBLocalization.KernelBranching.frameSign_orTangent_eq_one_iff},
\lean{Polyfold.SBLocalization.IsCanonicalFor.signedWeight_eq}}
\end{lemma}

\begin{lemma}[Positivity along paths]\label{lem:path}
Let $o$ be a continuous orientation of a tame sc-Fredholm section and let $q\mapsto(f+S_q)'(\gamma_q)$,
$q\in[0,1]$, be a family of isomorphisms along an sc-smooth path $\gamma$, with $S$ an sc$^+$-section of
the pulled-back bundle and $(f+S_q)(\gamma_q)=0$. Then $o$ is positive at $(f+S_0)'(\gamma_0)$ if and
only if it is positive at $(f+S_1)'(\gamma_1)$. The proof uses the local orientation propagation of
Theorems 6.6.11 and 6.6.13 (pp.~272--273).
\leanref{\lean{Polyfold.IsPosOr}, \lean{Polyfold.isPosOr_iff_of_path}}
\end{lemma}

\begin{lemma}[Invariance of continuous orientations]\label{lem:inv}
Let $o$ be a continuous orientation over an ep-groupoid whose objects need not all be smooth. Along a
continuous family of morphisms $\gamma\colon T\to X_1$ with $T$ connected and smooth sources and targets,
the condition that $o$ is invariant at $\gamma(t)$ does not depend on $t$.
\leanref{\lean{Polyfold.EpStrongBundle.isOrInvariantAt_iff_of_continuous}}
\end{lemma}

%======================================================================
\section{Remarks on scope}

\begin{enumerate}[label=(\arabic*)]
\item The first example is $\C$ with all levels equal, so its sc-structure is constant. The version of
  \S4 adds an infinite-dimensional factor with strictly nested levels, on which the section is the
  identity.
\item The orientation $o$ is chosen as the continuous orientation positive at the complex-linear
  isomorphisms. The library also contains an orientation \lean{sqOrientation}, one of the two
  continuous orientations, chosen without this normalization; its identification with $o$ is not part
  of the formal development, and the value of the count is stated for $o$. Theorems~\ref{thm:hyp}
  and~\ref{thm:admissible} are stated for both (\lean{canonicalGermInput_sqOrientationRepPos},
  \lean{exists_canonical_t5_pos}, \lean{count_eq_of_isAdmissible_pos}).
\item The description of $2/5$ as degree over isotropy order is not a formal statement, and the
  obstruction of Proposition~\ref{prop:obstruction} is stated for the finite-dimensional example only.
\end{enumerate}

%======================================================================
\appendix
\section{Statements and Lean declarations}

Files are in \lean{Polyfold/Examples/}; names written without a namespace are in
\lean{Polyfold.Examples.OrbifoldSquare}, and those of \S4 in \lean{Polyfold.Examples.OrbifoldSquareInf}.

{\footnotesize
\renewcommand{\arraystretch}{1.15}
\setlength{\tabcolsep}{4pt}
\begin{longtable}{@{}>{\raggedright\arraybackslash}p{0.22\textwidth}
                    >{\raggedright\arraybackslash}p{0.46\textwidth}
                    >{\raggedright\arraybackslash}p{0.27\textwidth}@{}}
\toprule
Statement & Lean declarations & File\\
\midrule
\endhead
\bottomrule
\endfoot
§1 $G\ltimes\C$, the weight-$2$ bundle, $z\mapsto z^2$
  & \lean{X5}, \lean{E5}, \lean{A5}, \lean{sqSection}
  & \lean{OrbifoldSquareBundle}\\
§1 effective; isotropy $\ZZ$ at $0$; the polyfold bundle and section
  & \lean{isEffective_X5}, \lean{isotropyZeroEquiv}, \lean{isRegular_X5_iff}, \lean{sqBundle},
    \lean{sqFredholmSection}
  & \lean{OrbifoldSquareStructure}\\
§1 sc-Fredholm, index $0$
  & \lean{isScFredholmSectionFunctor_sqSection}, \lean{isIndexAt_sqSection}
  & \lean{OrbifoldSquareOrientation}\\
§1 an orientation with tame germ data and index $0$
  & \lean{sqOrientationRep}, \lean{canonicalGermInput_sqOrientationRep}
  & \lean{OrbifoldSquareOrientation}\\
Prop.~\ref{prop:obstruction}, the linear algebra
  & \lean{linear_eq_zero_of_equivariant}
  & \lean{OrbifoldSquareAlgebra}\\
Prop.~\ref{prop:obstruction}
  & \lean{not_isTransversalAt_sectionMulti}
  & \lean{OrbifoldSquare}\\
§2 auxiliary norm; the multisection $\lambda_c$
  & \lean{isAuxNorm_n5}, \lean{controlsCompactness_n5}, \lean{Λ5}, \lean{t5},
    \lean{isScPlusMultisection_Λ5}, \lean{isInGeneralPosition_Λ5}, \lean{isControlledBy_t5}
  & \lean{OrbifoldSquarePerturbation}\\
Thm.~\ref{thm:hyp}
  & \lean{theorem_16_6_6_hypotheses}, \lean{theorem_16_6_6_sq}
  & \lean{OrbifoldSquare}\\
Thm.~\ref{thm:admissible}
  & \lean{exists_canonical_t5}, \lean{count_eq_of_isAdmissible}, \lean{exists_unique_count}
  & \lean{OrbifoldSquare}\\
Thm.~\ref{thm:count}, the two orbits
  & \lean{Λ5_sqSection_pos_iff_orbit}, \lean{π_ne_π_neg}, \lean{isRegular_of_Λ5_pos},
    \lean{Λ5_sqSection_eq}
  & \lean{OrbifoldSquare}\\
Thm.~\ref{thm:count}, the count as a sum over orbits; the form $1$
  & \lean{count_t5}, \lean{sqOne}, \lean{d_sqOne}
  & \lean{OrbifoldSquareCount}\\
Thm.~\ref{thm:count}, the orientation $o$ and the signs
  & \lean{sqOrientationPos}, \lean{sqOrientationRepPos}, \lean{positiveZeros_sqOrientationRepPos}
  & \lean{OrbifoldSquareSign}\\
Thm.~\ref{thm:count}
  & \lean{count_eq_two_fifths}, \lean{count_eq_two_fifths_all}
  & \lean{OrbifoldSquareSign}\\
§4 the weighted $\ell^2$ space
  & \lean{Polyfold.Examples.WeightedEllTwo.H},
    \lean{Polyfold.Examples.WeightedEllTwo.not_surjective_inclLevel},
    \lean{Polyfold.Examples.WeightedEllTwo.harmH_not_mem_smooth}
  & \lean{WeightedEllTwo}\\
Thm.~\ref{thm:inf}
  & \lean{theorem_16_6_6_hypotheses_inf}, \lean{theorem_16_6_6_inf},
    \lean{count_eq_two_fifths_inf}, \lean{count_eq_two_fifths_inf_all}
  & \lean{OrbifoldSquareInf}\\
Lemma~\ref{lem:finite}
  & \lean{Polyfold.EP.pointBranchedIntegral_eq_sum},
    \lean{Polyfold.SBLocalization.PointSolution.count_eq_sum}
  & \lean{PointCountFiniteSupport}\\
Lemma~\ref{lem:const}
  & \lean{Polyfold.MPolyfold.extD_constForm}
  & \lean{ConstantFormClosed}\\
Lemma~\ref{lem:sign}
  & \lean{Polyfold.SBLocalization.KernelBranching.frameSign_orTangent_eq_one_iff},
    \lean{Polyfold.SBLocalization.IsCanonicalFor.signedWeight_eq}
  & \lean{SignDimZeroFrames}\\
Lemma~\ref{lem:sign}, rank $0$: the continuation keeps the empty-wedge ray
  & \lean{Polyfold.SubspaceFamily.cont_emptyRay}
  & \lean{SignDimZeroLinear}\\
Lemma~\ref{lem:path}
  & \lean{Polyfold.IsPosOr}, \lean{Polyfold.isPosOr_iff_of_path}
  & \lean{SignDimZeroPush}\\
Lemma~\ref{lem:inv}
  & \lean{Polyfold.EpStrongBundle.isOrInvariantAt_iff_of_continuous}
  & \lean{OrInvariantSmooth}\\
\end{longtable}
}

\end{document}
